\section{Background and Motivation}

The rapid integration of large language models (LLMs) into everyday life has created a new category of first-responder for personal health concerns. Surveys conducted in the United States document that a substantial and growing proportion of people turn to general-purpose LLMs specifically for mental health support, often as a first contact before, or instead of, seeking professional care \citep{rousmaniere2025mentalhealth}. In parallel, a growing literature evaluates the clinical quality and safety of LLM-generated mental health responses in high-resource settings \citep{hua2025mentalhealth, santos2026clinical, saha2025linguistic}. What this literature has not addressed systematically is whether the quality and character of those responses remain stable across the geographic and infrastructural diversity of the users who now rely on them.

This gap is not merely academic. Mental health services are unevenly distributed across the world: in many low- and middle-income countries, specialist services are scarce or absent, community awareness of crisis support options is limited, and dedicated crisis helplines either do not exist or are not widely documented \citep{osborn2024students}. For a user in such a setting, an LLM may function not as a supplementary resource but as the primary available channel for mental health guidance. If the quality, actionability, or cultural framing of that guidance varies systematically with where the user says they are, the models most likely to be a user's only option are also the most likely to provide inferior support.

A second, less studied problem concerns the measurement of such variation. Much of the existing literature on LLM quality and bias in healthcare relies on automated content coding --- keyword indicators, dictionary-based feature extraction, or embedding-distance metrics --- to convert free-text model output into quantifiable outcomes at scale \citep{dai2023llm, depaoli2024thematic}. These pipelines are rarely validated against independent human judgements before their outputs are used to support substantive claims. This dissertation demonstrates directly, through one of its own measurement failures, that automated content-coding measures can produce outputs that are internally consistent and inter-rule-reliable yet still measure the opposite of what they claim to measure, with the error detectable only through comparison with independently collected human labels. This methodological observation is offered as a contribution to the LLM-audit literature in its own right.

\section{Problem Statement}

This dissertation addresses two related problems. The primary problem is empirical: whether the geographic location a user discloses when posing a mental health query to an LLM systematically changes the actionability, localisation, and cultural framing of the guidance the model provides, and whether any such variation runs in the direction that would disadvantage users in lower-income and lower-infrastructure settings. The secondary problem is methodological: whether automated content-coding pipelines used to measure such variation can be trusted without independent validation, and what a validation process that is capable of catching failures --- including directional failures --- looks like in practice.

\section{Research Aim and Objectives}

The aim is to determine whether and how the city and country stated by a user are associated with the mental-health guidance produced by contemporary large language models, while testing whether the automated measures used to quantify those differences are defensible.

The objectives are to construct a controlled factorial dataset for 20 cities; develop transparent text-based outcomes aligned with three pre-registered hypothesis families; compare those measures with human judgements using pre-specified criteria; test one statistical model per hypothesis while retaining city-level granularity; and assess sensitivity to the small number of geographic units, model exclusion, actionability components and contact verification. Measures redeveloped on the initial labels were frozen and evaluated on a fresh 160-response hold-out. Measures that failed this gate were retained only for exploratory hypothesis estimates.

\section{Research Questions}

Three research questions are pre-registered. \textbf{RQ1} asks whether actionability varies with the income category of the stated location. \textbf{RQ2} asks whether localisation varies with documented crisis-service availability. \textbf{RQ3} asks whether religious or spiritual recommendations vary by WHO region. A fourth, exploratory audit examines whether extracted crisis contacts are verified, visibly suspicious or unresolved; it is reported separately from the pre-registered hypothesis design. Validation gates determine whether each hypothesis result may be interpreted confirmatorily.

\section{Scope and Contribution}

The study covers 1,120 responses from seven models across 20 cities and eight mental-health scenarios. All queries are in English and no real user or patient data are used. The contribution is fourfold: sample-specific evidence for RQ1--RQ3; a documented case of automated content-coding failure; an exploratory audit of crisis-contact reliability; and a reproducible pipeline that records prompts, settings, rules, validation data and model outputs.

\section{Dissertation Structure}

Chapter 2 reviews relevant literature on LLMs in medical and mental health contexts, geographic and cultural bias in language model outputs, and existing approaches to LLM auditing. Chapter 3 describes the dataset construction, the NLP pipeline, the validation methodology, and the statistical models. Chapter 4 presents measurement validation before the hypothesis-model estimates and sensitivity checks. Chapter 5 discusses the findings in relation to the literature, their practical implications, and their limitations. Chapter 6 summarises the contributions and directions for future work.
